\documentclass{article}
\usepackage{graphicx} % Required for inserting images

\title{STT-2902}
\author{Julien Miron}


\begin{document}

\section*{Équipe 6 – Inférence à partir de données de sondage}

Objectif : Présenter comment les règles changent lorsqu’on travaille avec une population finie.


\vspace{1cm}
Suggestions :
\begin{itemize}
    \item Donnez un scénario concret (ex. : une classe, un village) et demandez aux autres d'estimer une proportion avec ou sans correction de population finie.
    \item Faites une comparaison visuelle des résultats (avec/sans correction).
    \item Expliquez dans quels cas cette correction est importante.

\end{itemize}

\vspace{1cm}
\textbf{Requis} :
\begin{enumerate}
    \item Montrer numériquement l’effet de la correction pour population finie sur un intervalle de confiance.
    \item Présenter un exemple où l’absence de correction mènerait à une conclusion erronée.
\end{enumerate}


\end{document}
